%% kapa: exam=1 round=19 year=2008 period=2 form=A subject=economics subject_ko=경제원론 seq=27 number=27
%% source: Q-Net 감정평가사 기출문제 (cst00309) artlSeq=1001962 file=extracted/1008/1001962/2008년 제19회 감정평가사 1차시험 기출문제/2교시.hwp
%% conversion: kordoc 4.24.51 -> tools/gen1.py (auto); review=proofread-sonnet; flags=none
\begin{question}{27}
두 재화만 소비하는 소비자의 소득소비곡선이 우하향한다. 이로부터 추론할 수 있는 것을 모두 고르면?
\begin{framedtext}
ㄱ. 두 재화가 보완재이다. \par
ㄴ. 두 재화가 모두 정상재이다. \par
ㄷ. 두 재화 중 한 재화만 열등재이다. \par
ㄹ. 두 재화 중 한 재화만 엥겔곡선이 우상향한다. \par
\end{framedtext}
\begin{choices}
  \item ㄱ, ㄴ
  \item ㄱ, ㄷ
  \item ㄴ, ㄹ
  \item ㄷ, ㄹ
  \item ㄱ, ㄷ, ㄹ
\end{choices}
\end{question}
\answer{4}
